% sec_metodologia.tex
% Sección 2: Materiales y Métodos
\section{Materiales y Métodos}

\subsection{Diseño Experimental}
Se diseñó un protocolo de validación estadística basado en simulaciones Monte Carlo con $K = 100$ ensayos independientes, cada uno con semilla pseudoaleatoria única generada mediante \texttt{std::random\_device}. Cada ensayo procesó $N = 5 \times 10^6$ operaciones de \textit{dispatch} condicional, totalizando $5 \times 10^8$ observaciones experimentales.

Se evaluaron cuatro estrategias de \textit{dispatch}:
\begin{itemize}
    \item \textbf{\texttt{switch} nativo:} Evaluación condicional en cada iteración (línea base).
    \item \textbf{\texttt{FastBranch}:} Propuesta experimental con punteros crudos y actualización infrecuente ($f_{\text{set}} = 1/1000$).
    \item \textbf{\texttt{FlexBranch}:} Variante con \texttt{std::function} para cuantificar el costo de la abstracción.
    \item \textbf{Control directo:} Llamada sin \textit{dispatch} condicional (límite teórico).
\end{itemize}

Las diferencias se validaron mediante un análisis estadístico no paramétrico de remuestreo (\textit{Bootstrapping}) con $B = 10,000$ iteraciones. Este enfoque permite construir intervalos de confianza (IC) empíricos al 95\%, mitigando las asunciones de normalidad paramétrica que suelen violarse en mediciones de latencia de hardware por interrupciones del SO. Se reporta adicionalmente el tamaño del efecto mediante la $d$ de Cohen.

\subsection{El Problema: Pipeline Flush}

La Fig.~\ref{fig:pipeline} ilustra el mecanismo por el cual una predicción fallida destruye el rendimiento del pipeline superescalar. Cuando el BPU falla, todas las instrucciones ejecutadas especulativamente se descartan, generando una burbuja de 14--18 ciclos.


% El diagrama de pipeline (Fig 1) ha sido trasladado a la discusión para ubicarse en la última página ocupando dos columnas.

\subsection{Implementación y Arquitectura}
Se desarrolló la biblioteca \texttt{semi\_static.hpp} (\textit{header-only}, C++17) con cuatro variantes (Cuadro~\ref{tab:clases}). La arquitectura central se basa en el desacoplamiento físico entre la evaluación de la condición (\textit{cold path}) y la ejecución del \textit{dispatch} (\textit{hot path}), como se ilustra en la Fig.~\ref{fig:arquitectura}.

\begin{table}[ht]
\centering
\caption{Variantes implementadas en \texttt{semi\_static.hpp}}
\label{tab:clases}
\begin{tabular}{@{}lccc@{}}
\toprule
\textbf{Clase} & \textbf{Mecanismo} & \textbf{Heap} & \textbf{Thread-safe} \\ \midrule
\texttt{FastBranch}       & Array \texttt{FuncPtr}   & No & No \\
\texttt{AtomicFastBranch} & Ídem + \texttt{atomic}   & No & Sí \\
\texttt{FlexBranch}       & \texttt{vector<function>} & Sí & No \\
\texttt{AtomicFlexBranch} & Ídem + \texttt{atomic}   & Sí & Sí \\ \bottomrule
\end{tabular}
\end{table}

\setcounter{figure}{1}
\begin{figure}[ht]
\centering
\begin{tikzpicture}[
    scale=0.85, every node/.style={scale=0.85}, % scaled down slightly to fit on 3 pages
    core/.style={rectangle, draw=gray!40, fill=gray!5, thick, rounded corners=2mm, minimum width=3.2cm, minimum height=0.8cm, align=center, font=\footnotesize\bfseries},
    cline/.style={rectangle, draw=gray!40, fill=gray!5, dashed, rounded corners=1mm, minimum width=3.8cm, minimum height=1.3cm},
    cell/.style={rectangle, draw=gray!50, fill=white, minimum width=1.4cm, minimum height=0.6cm, font=\scriptsize\ttfamily},
    arrow/.style={-{Stealth[scale=1.1]}, thick},
]

% Cache Line 1
\node[cline] (line1) at (-2.1, 0) {};
\node[font=\scriptsize\itshape, text=gray!80!black, anchor=south] at (line1.north) {Línea Caché A (Read-Only)};
\node[cell] (p0) at (-2.9, 0) {ptrs[0]};
\node[cell] (p1) at (-1.3, 0) {ptrs[1]};

% Cache Line 2
\node[cline, draw=techblue!50, fill=techblue!5] (line2) at (2.4, 0) {};
\node[font=\scriptsize\itshape, text=techblue, anchor=south] at (line2.north) {Línea Caché B (alignas 64)};
\node[cell, draw=techblue, thick] (idx) at (2.4, 0) {atomic\_idx};

% Threads
\node[core, draw=coldblue!60, text=coldblue] (core0) at (2.4, -2.2) {Cold Thread};
\node[core, draw=hotgreen!60, text=hotgreen!80!black] (core1) at (-0.5, 2.5) {Hot Thread(s)};

% Data flow
\draw[arrow, coldblue] (core0.north) -- (idx.south) node[midway, right, font=\scriptsize\ttfamily, text=coldblue] {store(release)};

\draw[arrow, hotgreen] (idx.north) -- (core1.south-|idx.north) node[midway, right, font=\scriptsize\ttfamily, text=hotgreen!80!black] {load(acquire)};

\draw[-{Stealth[scale=1.1]}, thick, hotgreen, dashed] (core1.south-|p1.north) -- (p1.north) node[midway, left, font=\scriptsize, text=hotgreen!80!black] {dispatch (call)};

% Annotation
\node[font=\scriptsize, text=alertred!80!black, align=center] at (0.1, -1.2) {\textbf{Aislamiento Físico}:\\Evita que el \texttt{store}\\invalide \texttt{ptrs}};

\end{tikzpicture}
\caption{Arquitectura de memoria y sincronización inter-hilo para HFT. La directiva \texttt{alignas(64)} aísla el índice atómico en una línea de caché exclusiva. Esto previene que las actualizaciones del \textit{Cold Thread} invaliden el caché L1 de los punteros a función (\textit{false sharing}). La semántica \textit{release/acquire} garantiza visibilidad de memoria sin utilizar costosos bloqueos.}
\label{fig:arquitectura}
\end{figure}

\subsection{Entorno de Ejecución}
Las pruebas se ejecutaron en un entorno Linux (Ubuntu 24.04 LTS), procesador Intel con 12 núcleos lógicos, compilador GCC~13.3 (\texttt{-O2}, C++17). La instrumentación temporal se realizó mediante \texttt{std::chrono::high\_resolution\_clock}. Las condiciones aleatorias se pre-generaron antes de cada \textit{trial} para asegurar que el costo del generador de números aleatorios (RNG) no contaminara las mediciones de latencia del \textit{hot path}.
